% ---------------------------------------------------------------------------
% OPTION A, CHAPTER 3. Assembled 2026-11-10 with AI assistance (Claude).
% Sections 3.5 and 3.6 (architecture, deployment) are the author's drafted
% text from the former design chapter (chapters/superseded/04-design.tex.txt),
% moved here with one cross-reference re-pointed. The other sections are new,
% written from docs/thesis/OUTLINE_BRIEF.md (chapter 3), CHAPTER_MATERIAL.md,
% BRIEF_VERIFICATION.md and docs/timeline/ on branch
% claude/thesis-brief-verification-6573e8. The use-case and domain diagrams
% (F4, F5) moved to appendix B.
% The author must read, revise and own every sentence before submission.
% ---------------------------------------------------------------------------
\chapter{Learnstation at a Glance}
\label{ch:system}

This chapter describes the system as a whole, so that the three chapters
that follow can each go deep into one part. It covers what the system is
for, the words it uses, what a reader can do in it today, how it is built and
deployed, and under which constraints. It answers no research question.

\section{From a Course Platform to a Reader}
\label{sec:vision}

Learnstation started in January 2026 as a platform for lecturers: a
professor creates a course, uploads slides, and students work through
lectures, quizzes and a dashboard. Its first idea was to stand between the
learning management system, which stores files, and the AI assistant, which
knows nothing about the course. Over the year the centre moved from the
course to the page. In the system that production serves since 29~October
2026, anyone drops a PDF and reads it at once, and everything else is added
to the page while it is read
(\code{docs/superpowers/specs/2026-10-26-library-and-spaces-design.md}).

The interface was therefore built twice. The second interface, called v4 in
the repository, was written as a new part of the frontend beside the first
one and not as a refactoring of it (\code{docs/design-v4/UI-SESSION-BRIEF.md},
PR~\#55). The repository records how this was done and not why refactoring
was ruled out. The author's reason is that the first interface assumed a
university, with professors and courses, and the second should work for any
material and any group. The older rooms still exist in the code and are no
longer where learners are sent. The backend changed as often:
\cref{sec:c4-lineage} lists about eight generations of the parsing pipeline
in ten months.

Five principles hold across the current system and are referred to in the
later chapters.
\begin{description}
  \item[Read first.] An uploaded file is readable before any AI step has
    finished. No feature may delay the reader.
  \item[The original page is the reference.] The page as the author made it
    is the default view. Generated text is one switch away and never
    replaces it.
  \item[Deterministic first, AI second.] What code can decide is decided by
    code: page hashes, the plan of a deck, the route of a page, the grading
    of a computable answer.
  \item[Parse once.] A file is processed once per uploader, and the result
    is reused for every reader of that upload.
  \item[The server decides.] Points, progress and grades are computed on the
    server; the browser displays them.
\end{description}

\section{Vocabulary and Roles}
\label{sec:stakeholders}

The system uses a small vocabulary on purpose
(\code{docs/V4\_GLOSSARY.md}). A \emph{Space} is a place where a group reads
the same material. A \emph{Lesson} is one uploaded document in a Space. A
\emph{page} is one page of that document, whether it was a slide or a sheet
of paper. An \emph{idea} is one thing a page teaches. The \emph{Library}
holds a person's own uploads, which the code calls materials. \emph{Luna} is
the assistant; she explains a page, asks about it, and runs short practice
sessions. The words course, lecture and slide belong to the older interface
and to the database, where a Space is still stored as a course
(\cref{sec:data-model}).

In the code a Space knows two roles: the \emph{owner}, who created it, and
the \emph{member}, who joined it
(\code{backend/services/v4\_spaces\_service.py:139--142}). Any signed-in
user can create a Space and can join a published one. Creating and
publishing is tied to ownership and not to an account type, so a student
can own a Space. A third role, an editor, appears in design documents and
is not built. The account types of the older interface (student, professor,
admin) still exist and still guard the older rooms; \cref{sec:authorisation}
shows one place where the two models met badly. \Cref{fig:use-case} in
\cref{app:figures} shows actors and goals as implemented, and
\cref{fig:domain-model} the vocabulary as a domain model.

\section{What Is Built and What Is Switched On}
\label{sec:built-vs-deployed}

A description of this system has to separate what exists in the code from
what a learner can reach. Features are guarded by twelve backend flags that
are all off by default, and production switches on four
(\cref{fig:flag-reality}): the v4 interface, its reader, sessions with Luna,
and uploads by students.

\thesisfigure{f17-flag-reality}
  {Built versus delivered. Twelve backend flags, all off by default;
   production switches on four. Since the end of October each interface flag
   follows its backend flag, which removed the half-states found at
   \thesisEarlierCommit{}.}
  {fig:flag-reality}

With these four flags a learner can upload a PDF, read it in two views (the
original page and a text written from it), ask Luna to explain a page, ask a
question about a page or about a whole Space and get an answer with links to
pages, answer checkpoint questions while reading, run a practice session
with Luna whose cards are scheduled by a variant of SM-2, take notes, and
see a history of earned points. An owner can publish a Space and, since two
days after the pinned revision, see how it is read (\cref{ch:loop}).

Several things are built and not reachable in the current interface: the
review engine of the older interface, an exam mode, generated study guides,
a global search page, and a third view of a page that rebuilds its layout.
We describe them only where a chapter needs them, and we do not count them
as delivered. Earlier in the year the flag of the interface and the flag of
the backend could disagree, so that a menu entry led to a route that
answered ``not found''. Two such mismatches existed. Since the end of
October the build derives each interface flag from its backend flag.

The thesis describes revision \code{\thesisCommit}. At the time of
writing (10~November 2026) production runs an earlier revision,
\code{c57721f2}, built on 6~November 2026, which lacks the last routing change of \cref{ch:content}
and the owner's view of \cref{ch:loop}. Production runs with all model
providers enabled; six have keys (probed on 9~November 2026).
\TODO{Author: re-check the deployed revision on the day of submission and
correct this paragraph.}

\section{Architecture}
\label{sec:architecture}

\thesisfigure{f6-components}
  {Component view of the deployed system. The api and the worker, built from
   one Dockerfile, call the LLM providers directly; the v4 app needs the direct
   asyncpg path to Postgres, which the readiness probe does not enforce.}
  {fig:components}

The system consists of three runtime parts and a set of managed services
(\cref{fig:components}). A single-page application written in React and
TypeScript runs in the browser. A FastAPI application~\cite{fastapi} serves the
HTTP API. A separate worker process, built on the Arq job queue~\cite{arq}, runs
everything that takes longer than a request should: parsing, LLM synthesis,
embedding, localisation, and scheduled maintenance. Supabase~\cite{supabase}
provides the PostgreSQL database, authentication and object storage. Two Redis
instances serve as cache and as job broker.

\paragraph{No gateway on the LLM path.}
The backend calls every LLM provider directly from one module
(\code{backend/services/ai/orchestrator.py}). Failover, rate limiting and cost
accounting live in that module, not in a separate service. This was not always
so. At \code{\thesisEarlierCommit} the production stack still contained a LiteLLM gateway
container, and the API carried its base URL in its environment, although no code
path sent traffic through it. The container was removed from the production
compose file on 28 August 2026, and the remaining configuration and legacy code
on 18 September 2026. The project's own documentation lagged behind: the backend
README still describes the gateway as the routing layer. Of all the
documentation drift discussed in \cref{sec:reality-vs-documented}, this case is
the most consequential, because it changes what design is being defended: not
``one gateway in front of many providers'', but two ordered provider chains
inside the application (\cref{sec:provider-strategy}).

\paragraph{Two database paths.}
The backend reaches PostgreSQL in two ways. The first is the Supabase client
over HTTP (PostgREST), in two variants: a service-role client that bypasses
row-level security, and an anonymous-key client that is subject to it. The
second is a direct connection pool (asyncpg) through Supabase's transaction
pooler. The pool is optional at start-up: if no database URL is configured, the
application starts without it; if one is configured but unreachable, start-up
fails fast. In practice the pool is not optional. The ingestion pipeline writes
its runs, lectures, slides and quiz questions through it, and every service of
the current (v4) application uses it. Older routes still run on the HTTP client;
the production compose file counts roughly 360 such call sites.

The two paths have a security consequence that shapes the rest of the design.
Writes through the pool and through the service-role client bypass row-level
security~\cite{postgresql-rls}, so for those routes authorisation is enforced in
application code, for example by an explicit course-ownership check before an
upload is accepted. Row-level security is the boundary only for queries the
browser makes directly. \Cref{sec:authorisation} returns to this.

\paragraph{The readiness signal.}
The API exposes a liveness probe (\code{/health}) and a readiness probe
(\code{/health/ready}). The readiness probe checks the database, both Redis
instances and whether a worker is alive when jobs are waiting. One condition
deserves mention: when no database URL is configured, the probe records the
database as \code{not\_configured} and still reports ready. Routes that use the
Supabase HTTP client still work without the pool, but the pipeline and the v4
routes do not, so this is not a usable mode; at \code{\thesisPreviousCommit} the
code comment still called it a supported ``REST-only'' mode, and it now says
plainly that every upload fails while the check reports ready. In
production the URL is set and the pool is initialised at start-up, so this is a
latent gap rather than a live fault: a misconfiguration would pass the health
check and surface only as failing requests.

\paragraph{Browser to backend.}
The browser talks to Supabase directly only for authentication, for reading the
user's role, and in a few remaining screens of the older interface. All data of
the v4 application goes through the FastAPI backend, which receives the user's
Supabase session token as a bearer token. A test enforces this boundary: no
module of the v4 feature folder may import the Supabase client.

\section{Deployment}
\label{sec:deployment}

\thesisfigure{f7-deployment}
  {Production deployment as run since 2026-10-30: five containers on one
   Hetzner VM behind a host nginx that terminates TLS, with Supabase, Resend,
   Better Stack and Cloudflare (DNS, R2) as external services. A systemd timer
   pings Better Stack every five minutes; the nightly dump is age-encrypted
   before upload to R2, and its decryption key is kept off the server.}
  {fig:deployment}

Production runs on one virtual machine (Hetzner CX23: two virtual CPUs, 4\,GB of
memory, Ubuntu~24.04) with five containers defined in
\code{docker-compose.prod.yml} (\cref{fig:deployment}):
\begin{itemize}
  \item \code{frontend}: nginx serving the built application on port 3000, bound
    to the loopback interface only, and forwarding \code{/api/} to the API;
  \item \code{api}: FastAPI under one uvicorn worker, limited to 512\,MB;
  \item \code{worker}: the Arq worker, also limited to 512\,MB, running up to four
    jobs at once;
  \item \code{redis}: the cache, 256\,MB with least-recently-used eviction;
  \item \code{redis-queue}: the job broker, 128\,MB with eviction disabled and
    append-only persistence, because a queue that silently drops jobs is worse
    than one that refuses new ones.
\end{itemize}
TLS is terminated by nginx on the host, with certificates from Let's Encrypt.
That host configuration lives on the server, not in the repository; the deploy
runbook documents it. DNS is held by Cloudflare without proxying. The database,
authentication and storage are a managed Supabase project in the EU (Ireland)
region on the free tier, which provides no point-in-time recovery.

\paragraph{A reproducibility finding.}
Until November the repository held four compose files: one for local
development, a production file whose header named ``a shared university
server'', a staging file that referred to Redis hosts it did not define, and one
for Oracle Cloud. None named the host actually used. The staging and Oracle files
were deleted and the production header now names the Hetzner machine (found,
then fixed). What still exists only on the server is the host proxy
configuration and the nightly database dump script, so the deployed system
still cannot be reconstructed from the repository alone. The gap is narrower
than it was in August: the deploy, monitoring and backup procedures are written
down as runbooks, and the heartbeat and off-site backup scripts are versioned.
But the
system's ground truth is still the server, not the repository; this is why the
project's own rules say to inspect the running images rather than trust
\code{git log}.

\paragraph{Operations.}
Deploys are manual: a build and restart over SSH, with the previous images
tagged for rollback. A systemd timer runs a heartbeat every five minutes. It
checks readiness, that all five containers run, that the public site answers,
disk and memory headroom, and certificate validity, and it reports to an
external monitor that alerts when the heartbeat goes silent. The nightly
database dump is encrypted with \code{age} and copied to object storage in the
EU, with a 14-day retention. These are modest measures, sized to a system with
very few real users, which is the scale the platform actually operates at.

\section{Technology Choices}
\label{sec:stack}

Most of the stack was not chosen by deliberation. React, Vite, TypeScript,
Tailwind and Supabase came with the scaffold of the first commit in January
(\code{184e5742}), and FastAPI, Redis and nginx followed as the common
companions. We name them as defaults and do not defend them. Seven choices
do have a recorded reason, and they are listed in \cref{tab:decisions}.

\begin{table}[htbp]
  \centering\small
  \caption{Technology decisions with a recorded reason (dates are 2026).}
  \label{tab:decisions}
  \begin{tabular}{@{}p{0.26\textwidth}p{0.46\textwidth}p{0.20\textwidth}@{}}
    \toprule
    \textbf{Decision} & \textbf{Recorded reason} & \textbf{Source} \\
    \midrule
    A background job queue (Arq) & Parsing takes longer than a request may. &
      \code{6ce7b5f1}, 19~May \\
    Two Redis instances & A cache may evict; a queue must not lose a job. &
      \code{43a86927}, 16~July \\
    PyMuPDF for text and rendering & A second renderer was dropped for a
      memory leak; the plain text call stayed as the only live extractor. &
      \code{227267e7}, \code{b1c7a428} \\
    No LLM gateway & The gateway had no live callers, used 1\,GB of a 4\,GB
      host, and its health check could block start-up. &
      \code{68a091e0}, 25~Aug \\
    HNSW instead of IVFFlat & Recorded in the migration that replaces the
      index. & \code{20260719000002} \\
    Small WebP posters for page previews & Full PDFs as backdrops exceeded
      the storage egress quota; posters cut egress by 97.6\,\%. &
      PR~\#12, 17~Aug \\
    One embedding size (768) & Two embedding models fill the same column. &
      \code{ai/embeddings.py} \\
    \bottomrule
  \end{tabular}
\end{table}

The third row hides a finding that \cref{ch:content} measures: the plain
text call became the production extractor by attrition and not by
comparison, and the benchmark places it second to last.

\section{Constraints}
\label{sec:constraints}

Four constraints shaped the design more than any requirement. The system
runs on one virtual machine with two virtual CPUs and 4\,GB of memory.
Model calls run on free or low-cost quotas, which is why every chain of
providers in \cref{ch:content} has a fallback and a spending cap. The
database is a managed PostgreSQL on a free tier without point-in-time
recovery. And the team was two students, then one.

One more condition is not a constraint and still shapes every evaluation in
this thesis: the system has almost no users. The evaluations therefore work
on material and not on people. The material is 25~lecture decks with
1{,}593~pages from two courses of one university, in English and German,
plus one further deck and two exercise sheets for the benchmark
(\cref{sec:c4-requirements}).
